\centerline{\bf \Huge Зимняя Сессия}
\centerline{\bf \Huge }

\begin{flushright}
        \scriptsize{\textbf{14:31. Идёт снег. Прибыла Елена Михайловна. Начало Зимней Сессии. Выживут сильнейшие}}
\end{flushright}

\problem{1} Решите уравнения:

а) $|2X-30|=6$

б) $\dfrac{8}{15}Y - \dfrac{3}{10}Y + \dfrac{3}{5}Y = \dfrac{1}{2}$

в) $X : 5\dfrac{5}{8} = 6,4 : 7\dfrac{1}{2}$

\problem{2} Теплоход прошел за 0,75 часа по течению реки 16,62 км, а против течения за 0,8 часа 14,4 км. Какова скорость теплохода в стоячей воде, скорость реки, скорость теплохода по течению, скорость теплохода против течению?

\problem{3} Половину пути мотоциклист ехал со скоростью 45 км/ч, затем задержался на 10 минут и чтобы наверстать время, он увеличил скорость на 15 км/ч. Каков весть путь?

\problem{4} Елена Михайловна съедает новогодний подарок за 6 дней, Анджур Аркадьевич за 3 дня, а Расул Владимирович за 2 дня. За сколько дней съедят новогодний подарок вместе все учителя ЭлМШ?